\subsection{Mathematical Research and Integrity}

Mathematical research alternates discovery, attempted refutation, proof, verification, and communication rather than proceeding linearly from question to theorem.

\begin{process} \label{proc:mathematical-research-cycle} \hlt{Mathematical Research Cycle}\\
For a proposed research problem, maintain a revisable cycle with the following stages.
\begin{enumerate}[label=(\roman*)]
\setlength{\itemsep}{0pt}
\item \hlt{Question}: state the object, scope, known assumptions, and a falsifiable or provable target.
\item \hlt{Literature}: identify definitions, known results, counterexamples, adjacent formulations, and open status.
\item \hlt{Exploration}: compute examples, boundary cases, and nonexamples while recording the tested domain.
\item \hlt{Conjecture}: state the strongest precise claim supported by exploration, with all suspected hypotheses.
\item \hlt{Attack}: search independently for a proof and a counterexample; record failed routes and new obstructions.
\item \hlt{Revision}: weaken, strengthen, split, or reframe the conjecture when an example or proof obligation demands it.
\item \hlt{Verification}: check every proof dependency, compare with prior results, and seek independent criticism.
\item \hlt{Communication}: state novelty accurately, write a complete proof, cite provenance, and record limitations.
\end{enumerate}
MIT's project laboratory similarly presents mathematical discovery as data generation, conjecture, experiment, proof or counterexample, and writing \citep{mit2013projectlaboratory}.
\end{process}

\begin{process} \label{proc:research-claim-ledger} \hlt{Research Claim Ledger}\\
Maintain one row per substantive claim with its exact wording, status, evidence, dependencies, source locator, and unresolved checks. Distinguish proved, conjectured, computed, observed, reproduced, replicated, and interpreted claims. Update the row when a hypothesis, dataset, code revision, or conclusion changes. This record prevents an exploratory observation from silently becoming a theorem or empirical conclusion.
\end{process}

\begin{convention} \label{conv:reproducibility-replicability-generalizability} \hlt{Reproducibility, Replicability, and Generalizability}\\
Use the National Academies terminology \citep{nasem2019reproducibility}.
\begin{enumerate}[label=(\roman*)]
\setlength{\itemsep}{0pt}
\item \hlt{Reproducibility}: consistent computational results from the same data, code, methods, and analysis conditions.
\item \hlt{Replicability}: consistent results from a new study addressing the same question with newly obtained data.
\item \hlt{Generalizability}: applicability of a result to contexts or populations beyond those directly studied.
\end{enumerate}
These are separate targets. A reproducible computation need not replicate on new data or generalize to a new setting.
\end{convention}

\begin{principle} \label{princ:research-integrity-record} \hlt{Research Integrity and Stewardship}\\
Research records should permit scrutiny, correction, and appropriate reuse. Report methods and findings honestly; preserve accurate records; credit contributions and prior work; disclose conflicts; respect confidentiality, consent, licenses, and governing rules; and correct the record when an error is found. These duties reflect the honesty, accountability, fairness, and stewardship principles of the Singapore Statement \citep{singapore2010integrity}.
\end{principle}

\begin{warning} \label{warn:novelty-and-negative-results} \hlt{Novelty and Negative Results}\\
Do not equate novelty with truth or publication with correctness. Search for prior and concurrent work before claiming priority, and report a failed proof, null result, or nonreplication with its actual protocol and uncertainty. A negative outcome may diagnose a false conjecture, an inadequate method, or a boundary of generality; the evidence must distinguish these possibilities.
\end{warning}
